\chapter{Vision impairment}
\index{Vision impairment}

\section*{Definition and Overview}
Vision impairment is defined as a functional limitation of the visual system characterized by a decrease in visual acuity or a restriction in the visual field that cannot be corrected to normal levels by conventional means, such as spectacles, contact lenses, medication, or surgical interventions. The World Health Organization (WHO) classifies vision impairment into distance and near vision impairment. Distance vision impairment is categorized based on visual acuity in the better-seeing eye with best possible correction: mild (visual acuity worse than 6/12 to 6/18), moderate (visual acuity worse than 6/18 to 6/60), severe (visual acuity worse than 6/60 to 3/60), and blindness (visual acuity worse than 3/60 or a corresponding visual field loss of less than 10 degrees around central fixation). Near vision impairment is defined as a presenting near visual acuity worse than N6 or M.08 at a distance of 40 centimetres.

The significance of vision impairment extends far beyond the clinical domain; it is a profound public health, economic, and social challenge. Vision is the primary sensory modality through which humans interact with their environment, accounting for the vast majority of sensory input. Consequently, any significant degradation in visual function impairs an individual's ability to perform activities of daily living, such as reading, writing, cooking, and navigating physical spaces. This limitation leads to a loss of personal autonomy, restricted mobility, and a reduced capacity for employment, which can precipitate economic hardship. Furthermore, vision impairment is closely linked to adverse mental health outcomes, including clinical depression, anxiety, and social isolation, particularly in elderly populations. By understanding the multi-faceted nature of vision impairment, healthcare providers and policymakers can better design targeted interventions to mitigate its global burden.

\section*{Epidemiology}
The epidemiological landscape of vision impairment is vast and characterized by significant demographic and geographic disparities. Globally, the WHO estimates that at least 2.2 billion people have a near or distance vision impairment. Of these cases, approximately 1 billion involve vision impairment that could have been prevented or has yet to be addressed. The leading causes of distance vision impairment and blindness on a global scale are uncorrected refractive errors and cataracts. Other major contributors include age-related macular degeneration, glaucoma, diabetic retinopathy, corneal opacification, and trachoma.

Demographic analysis reveals that the prevalence of vision impairment is not uniformly distributed. It is heavily skewed toward older age groups, with the majority of affected individuals being aged 50 years or older. As global life expectancy increases and populations age, the absolute number of individuals experiencing age-related visual decline is projected to rise exponentially. Gender disparities also exist, with women experiencing higher rates of vision impairment and blindness than men across most regions. This discrepancy is attributed to a combination of longer female life expectancy and systematic barriers in accessing eye care services in various parts of the world.

Geographically, low- and middle-income countries bear a disproportionate share of the global burden. More than 89\% of visually impaired individuals reside in developing nations, particularly in sub-Saharan Africa, South Asia, and parts of East Asia. In these regions, the lack of ophthalmological infrastructure, shortage of trained eye care professionals, and high costs of treatment prevent timely diagnosis and intervention. Consequently, diseases that are easily treatable or preventable in high-income nations, such as senile cataracts or refractive errors, remain major causes of avoidable blindness in poorer regions.

\section*{Aetiology and Pathophysiology}
The development of vision impairment involves diverse anatomical structures and pathophysiological mechanisms. The principal causes and their mechanisms include:
\begin{itemize}
    \item \textbf{Uncorrected Refractive Errors}: These arise when the optical system of the eye (the cornea and crystalline lens) fails to focus incident light rays precisely onto the light-sensitive photoreceptor layer of the retina. In myopia (nearsightedness), the axial length of the globe is too long relative to the refractive power, or the cornea is excessively curved, causing images to focus in front of the retina. Conversely, in hyperopia (farsightedness), the axial length is too short or the corneal curvature is flat, causing the focal point to lie behind the retina. Astigmatism occurs due to an irregular, asymmetric curvature of the cornea or lens, leading to multiple focal points and distorted vision. Presbyopia is an age-related loss of accommodation resulting from the progressive hardening of the crystalline lens and decrease in ciliary muscle elasticity, which limits the eye's ability to focus on near objects.
    \item \textbf{Cataract}: A cataract is characterized by the progressive opacification of the crystalline lens. The lens is composed of highly ordered water and protein molecules called crystallins, which maintain optical clarity. With age, exposure to ultraviolet (UV) radiation, oxidative stress, trauma, or metabolic disorders (such as diabetes), these proteins undergo denaturing, aggregation, and pigmentation. This disruption of the organized micro-architecture scatters and absorbs incident light, preventing it from reaching the retina clearly, and manifests as a slow, painless decline in visual acuity and contrast sensitivity.
    \item \textbf{Glaucoma}: Glaucoma represents a group of progressive optic neuropathies characterized by the degeneration of retinal ganglion cells (RGCs) and their axons, leading to a distinctive cupping of the optic nerve head and corresponding visual field deficits. The primary modifiable risk factor is elevated intraocular pressure (IOP), which is typically caused by impaired drainage of aqueous humour through the trabecular meshwork and Schlemm's canal. The pathophysiology involves mechanical compression of the optic nerve axons at the lamina cribrosa, microvascular ischaemia of the optic nerve head, and neurotoxicity mediated by glutamate and oxidative stress, leading to apoptotic cell death of the RGCs.
    \item \textbf{Age-Related Macular Degeneration (AMD)}: AMD targets the macula lutea, the central region of the retina responsible for high-acuity, detailed vision. It is divided into dry (non-neovascular) and wet (neovascular) forms. Dry AMD is characterized by the accumulation of extracellular metabolic debris, known as drusen, between the retinal pigment epithelium (RPE) and Bruch's membrane. Over time, this leads to progressive atrophy of the RPE and overlying photoreceptors (geographic atrophy). Wet AMD is characterized by the pathological growth of abnormal choroidal blood vessels (choroidal neovascularization) across Bruch's membrane, stimulated by vascular endothelial growth factor (VEGF). These fragile vessels leak fluid, lipids, and blood into the subretinal space, causing rapid, severe, and often irreversible central vision loss.
    \item \textbf{Diabetic Retinopathy}: This is a microvascular complication of chronic hyperglycaemia. Prolonged high blood glucose levels induce biochemical cascades (including the polyol pathway and advanced glycation end-product accumulation) that damage retinal capillary endothelial cells and pericytes. This leads to capillary basement membrane thickening, microaneurysm formation, and increased vascular permeability, resulting in macular oedema (non-proliferative diabetic retinopathy). As capillary occlusion progresses, retinal ischaemia triggers the upregulation of angiogenic factors, predominantly VEGF. This stimulates the proliferation of fragile new blood vessels on the retinal surface and into the vitreous cavity (proliferative diabetic retinopathy), which are prone to vitreous haemorrhage and tractional retinal detachment.
    \item \textbf{Trachoma}: An infectious disease caused by specific serotypes (A, B, Ba, and C) of the bacterium \textit{Chlamydia trachomatis}. Transmitted via direct contact or eye-seeking flies, repeated reinfection in childhood leads to chronic follicular conjunctivitis. The subsequent inflammatory response causes scarring of the tarsal conjunctiva, which deforms the eyelid, pulling the eyelashes inward (entropion and trichiasis). The constant friction of the eyelashes against the cornea leads to corneal ulceration, secondary bacterial infections, and eventual permanent corneal opacification and blindness.
\end{itemize}

\section*{Clinical Features}
The clinical presentation of vision impairment depends on the underlying pathological process, the anatomical structures involved, and the stage of the disease. Key signs, symptoms, and clinical features include:
\begin{itemize}
    \item \textbf{Reduction in Visual Acuity}: Patients typically report blurred vision at distance, near, or both. This may manifest as difficulty reading fine print, recognizing familiar faces, reading street signs, or watching television. In refractive errors, this blur is present across all distances or specific ranges, whereas in central macular pathology, it specifically affects the central visual field.
    \item \textbf{Visual Field Defects}: Impairment can affect the peripheral or central vision. Peripheral visual field loss, often referred to as ``tunnel vision,'' is a classic feature of advanced glaucoma and retinitis pigmentosa. Central scotomas (blind spots in the center of vision) are indicative of macular disease, such as dry or wet AMD, or optic nerve pathology.
    \item \textbf{Metamorphopsia}: The distortion of visual images, where straight lines appear wavy, bent, or distorted. This is a hallmark symptom of macular oedema, epiretinal membranes, and active choroidal neovascularization in wet AMD. It can be clinically assessed using an Amsler grid.
    \item \textbf{Reduced Contrast Sensitivity and Colour Vision Deficits}: Affected individuals often struggle to distinguish objects of similar tone or brightness, making navigation in low-light environments difficult. Impaired colour vision (dyschromatopsia) can occur, where colours appear washed out, faded, or difficult to differentiate, commonly seen in optic neuropathies and advanced cataracts.
    \item \textbf{Nyctalopia (Night Blindness)}: A significant difficulty or inability to see in dim light or dark environments. This is a common feature of rod-photoreceptor dysfunction, vitamin A deficiency, and advanced glaucoma.
    \item \textbf{Photophobia and Glare Sensitivity}: Extreme sensitivity to bright lights, sunlight, or headlights. Cataract patients frequently complain of halos and glare around light sources, which is caused by light scattering within the opacified lens.
    \item \textbf{Ocular Asthenopia}: Eye strain, physical fatigue, and headaches associated with prolonged visual tasks. This is highly prevalent in patients with uncorrected refractive errors or binocular vision anomalies.
\end{itemize}

\section*{Differential Diagnosis}
When evaluating a patient presenting with vision impairment, it is essential to differentiate between ocular, systemic, and neurological conditions that can present with overlapping visual symptoms. The differential diagnoses to consider include:
\begin{itemize}
    \item \textbf{Stroke and Transient Ischaemic Attack (TIA)}: Acute cerebrovascular accidents affecting the visual pathways (e.g., the optic radiations or visual cortex in the occipital lobe) can cause sudden, painless visual field loss, such as homonymous hemianopia. A TIA may present as transient monocular blindness (amaurosis fugax) due to microemboli temporarily occluding the ophthalmic or central retinal artery.
    \item \textbf{Optic Neuritis}: An inflammatory, demyelinating condition of the optic nerve, frequently associated with multiple sclerosis. It presents with acute, typically unilateral, painful vision loss, impaired colour vision, and a relative afferent pupillary defect (RAPD). The pain is characteristically exacerbated by eye movements.
    \item \textbf{Retinal Detachment}: The physical separation of the neurosensory retina from the underlying retinal pigment epithelium. Symptoms include the sudden onset of photopsias (flashes of light), a sudden increase in vitreous floaters, and a progressive dark shadow or ``curtain'' obscuring a portion of the visual field. It is a sight-threatening emergency.
    \item \textbf{Giant Cell Arteritis (Temporal Arteritis)}: A systemic autoimmune vasculitis of medium- and large-sized arteries that represents a medical emergency. If it involves the posterior ciliary arteries, it causes arteritic anterior ischaemic optic neuropathy (AAION), leading to sudden, profound, and permanent visual loss. Associated symptoms include jaw claudication, scalp tenderness, temple headache, fever, and polymyalgia rheumatica.
    \item \textbf{Keratoconus}: A progressive, non-inflammatory thinning and ectasia of the cornea, causing it to assume a conical shape. This results in high irregular astigmatism and progressive visual impairment that cannot be adequately corrected with standard spectacles, requiring rigid gas-permeable contact lenses or corneal collagen cross-linking.
    \item \textbf{Retinal Vascular Occlusions}: Central or branch retinal artery occlusions (CRAO/BRAO) cause sudden, painless, severe visual loss, often characterized by a ``cherry-red spot'' on the macula on fundoscopy. Central or branch retinal vein occlusions (CRVO/BRVO) present with variable subacute or acute visual loss and characteristic ``blood and thunder'' retinal haemorrhages.
    \item \textbf{Functional or Psychogenic Visual Loss}: Non-organic visual impairment where there is no anatomical or physiological pathology to explain the visual deficits. This includes conversion disorder and malingering, and is diagnosed through specialized testing (e.g., visual evoked potentials or incongruous visual field patterns).
\end{itemize}

\section*{When to Seek Medical Care}
Any change in vision warrants professional evaluation, but certain symptoms indicate acute, potentially sight- or life-threatening emergencies that require immediate medical attention. Individuals must seek emergency care if they experience:
\begin{itemize}
    \item Sudden, rapid, or complete loss of vision in one or both eyes, even if the vision returns to normal shortly thereafter.
    \item Sudden onset of severe, deep eye pain, which may be accompanied by physical redness, headache, nausea, vomiting, or halos around lights (suspicious for acute angle-closure glaucoma).
    \item The sudden appearance of flashes of light, a shower of dark floaters, or a grey or black shadow moving across the field of vision.
    \item Physical trauma, penetrating injury, or chemical exposure to the eye or surrounding facial structures.
    \item Double vision (diplopia) that is sudden in onset.
    \item Vision loss accompanied by systemic neurological signs, such as weakness or numbness of the face or limbs, difficulty speaking, confusion, loss of balance, or difficulty walking.
\end{itemize}
For any of these acute or life-threatening symptoms, patients should immediately contact their local emergency services (e.g., local emergency services in many countries, 911 in the United States, 111 or 999 in the United Kingdom) or present to the nearest hospital emergency department with specialized ophthalmic services.

\section*{Diagnosis and Investigations}
Confirming the diagnosis, identifying the specific cause, and staging the severity of vision impairment requires a systematic clinical assessment and specialized diagnostic investigations:
\begin{itemize}
    \item \textbf{Visual Acuity and Refraction}: Standardized measurement of visual acuity is performed using Snellen or LogMAR charts at distance (typically 6 metres or 20 feet) and near (40 centimetres). Subjective and objective refraction (using an autorefractor or retinoscope) determines the presence and magnitude of refractive errors and the best-corrected visual acuity (BCVA).
    \item \textbf{Slit-Lamp Biomicroscopy}: High-magnification examination of the anterior segment of the eye, including the eyelids, conjunctiva, cornea, anterior chamber, iris, and crystalline lens. This allows for the detection of corneal scars, inflammatory cells in the aqueous humour, iris abnormalities, and cataracts.
    \item \textbf{Tonometry}: Measuring intraocular pressure (IOP) is critical for glaucoma screening and monitoring. Goldmann applanation tonometry remains the clinical gold standard, though air-puff tonometry and rebound tonometry are also utilized.
    \item \textbf{Gonioscopy}: The use of a specialized mirrored lens on the cornea to visualize the iridocorneal angle, allowing the clinician to differentiate between open-angle and closed-angle glaucoma.
    \item \textbf{Fundoscopy (Ophthalmoscopy)}: Direct or indirect evaluation of the posterior segment of the eye under pharmacologically induced pupillary dilation (mydriasis). The clinician inspects the optic nerve head (for cup-to-disc ratio and pallor), the macula (for drusen, haemorrhage, or fluid), the retinal vasculature (for diabetic or hypertensive changes), and the peripheral retina (for tears or detachments).
    \item \textbf{Optical Coherence Tomography (OCT)}: A non-invasive, high-resolution optical imaging technique that provides cross-sectional micron-scale images of the retinal layers. It is crucial for diagnosing and monitoring macular diseases, such as macular oedema, wet AMD, macular holes, and for measuring the retinal nerve fibre layer (RNFL) thickness in glaucoma.
    \item \textbf{Visual Field Perimetry}: Automated perimetry (such as the Humphrey Visual Field analyzer) systematically maps the patient's visual sensitivity to light points across the visual field. This is vital for detecting and tracking the peripheral vision loss characteristic of glaucoma and neurological pathway lesions.
    \item \textbf{Fluorescein Angiography}: A diagnostic test involving the injection of a fluorescent dye into a systemic vein, followed by rapid sequential photography of the retinal circulation. This identifies areas of capillary non-perfusion, microaneurysms, and active dye leakage from abnormal neovascular membranes.
\end{itemize}

\section*{Management}
The management of vision impairment is highly individualized and focuses on correcting the refractive state, treating the underlying disease process, and rehabilitating the patient to maximize functional independence:
\begin{itemize}
    \item \textbf{Refractive and Optical Corrections}: Uncorrected refractive errors are managed with prescription spectacles or contact lenses (soft, rigid gas-permeable, or scleral). Surgical options include corneal refractive surgeries such as laser-assisted in situ keratomileusis (LASIK), photoreactive keratectomy (PRK), or the implantation of phakic intraocular lenses.
    \item \textbf{Surgical Treatments}: Cataract surgery is one of the most common and successful surgical procedures performed globally. It involves removing the cloudy natural lens via phacoemulsification and replacing it with a clear synthetic intraocular lens (IOL). For glaucoma, surgeries such as trabeculectomy, glaucoma drainage devices (aqueous shunts), or minimally invasive glaucoma surgery (MIGS) are performed to lower IOP. Vitreoretinal surgery (vitrectomy) is indicated for non-clearing vitreous haemorrhage, tractional retinal detachment, and macular holes.
    \item \textbf{Pharmacotherapy}: Glaucoma is primarily managed with topical ophthalmic drops that reduce aqueous humour production (e.g., beta-blockers, alpha-agonists, carbonic anhydrase inhibitors) or increase its outflow (e.g., prostaglandin analogues, rho-kinase inhibitors). For wet AMD, diabetic macular oedema, and retinal vein occlusions, first-line therapy involves regular intravitreal injections of anti-VEGF agents (such as aflibercept, ranibizumab, or faricimab) to suppress angiogenesis and reduce vascular leakage.
    \item \textbf{Laser Therapy}: Various ophthalmic lasers are used for target-specific treatments: Selective Laser Trabeculoplasty (SLT) to improve outflow in open-angle glaucoma; Laser Peripheral Iridotomy (LPI) to relieve pupillary block in angle-closure glaucoma; and Panretinal Photocoagulation (PRP) to ablate ischaemic peripheral retina in proliferative diabetic retinopathy, reducing the drive for neovascularization.
    \item \textbf{Low Vision Rehabilitation and Assistive Technologies}: For individuals with irreversible vision loss, low vision clinics provide specialized support. This includes prescribing optical magnifiers, electronic desktop magnifiers, closed-circuit television (CCTV) systems, and telescopic aids. Adaptive technologies, such as screen-reading software (e.g., JAWS or NVDA), text-to-speech programs, voice-controlled smart home devices, and mobile applications designed for the visually impaired, significantly enhance independence. Orientation and mobility training teaches patients how to navigate environments safely using long canes or guide dogs.
\end{itemize}

\section*{Complications}
Untreated, progressive vision impairment or its interventions can lead to significant physical, psychological, and systemic complications:
\begin{itemize}
    \item \textbf{Irreversible Blindness}: Chronic conditions such as advanced glaucoma, diabetic retinopathy, and wet AMD, if left untreated, progress to complete and permanent destruction of the visual pathway and total blindness.
    \item \textbf{Increased Incidence of Falls and Fractures}: Visual deficits, particularly those affecting depth perception, contrast sensitivity, and peripheral vision, impair spatial orientation and increase the risk of falls, leading to serious injuries such as hip fractures, head trauma, and soft tissue damage.
    \item \textbf{Psychological Morbidity}: The loss of independence, inability to engage in occupational and recreational activities, and difficulty navigating social interactions lead to high rates of clinical depression, generalized anxiety, and social withdrawal.
    \item \textbf{Accelerated Cognitive Decline}: Longitudinal studies demonstrate a significant association between untreated vision impairment and an accelerated rate of cognitive deterioration and dementia in older adults, likely due to reduced sensory input and subsequent social and mental disengagement.
    \item \textbf{Treatment-Related Complications}: Medical and surgical interventions carry inherent risks. Intravitreal injections can lead to endophthalmitis (a severe, sight-threatening intraocular infection), intraocular pressure spikes, or retinal tears. Intraocular surgery can be complicated by posterior capsule rupture, corneal decompensation, macular oedema, or choroidal haemorrhage.
    \item \textbf{Socioeconomic Deprivation}: Visual impairment can limit employment opportunities, decrease productivity, and necessitate costly continuous medical care, imposing a severe financial strain on patients, their families, and the broader healthcare system.
\end{itemize}

\section*{Prognosis and Outcomes}
The prognosis and long-term outcomes for individuals with vision impairment are highly variable and depend primarily on the underlying cause, the stage of the disease at the time of diagnosis, and the patient's adherence to treatment. For conditions like uncorrected refractive errors and uncomplicated cataracts, the prognosis is exceptionally positive. Modern surgical and optical interventions can restore normal or near-normal vision in the vast majority of cases, resulting in immediate improvements in quality of life.

For progressive, chronic diseases such as glaucoma, dry AMD, and diabetic retinopathy, the prognosis is more guarded. Because damaged retinal ganglion cells and photoreceptors cannot regenerate, any vision loss that has occurred prior to treatment is irreversible. The therapeutic goal in these cases is not to restore vision, but to preserve remaining visual function and slow the rate of disease progression. When these conditions are detected early through routine screening and managed proactively with high compliance to therapy, patients can often maintain functional vision for the remainder of their lives. Conversely, late-stage presentation or poor compliance leads to a significantly worse prognosis, often resulting in severe visual disability or total blindness.

\section*{Prevention and Self-Care}
Preventative strategies and proactive self-care are essential to reducing the global incidence of vision impairment and preserving ocular health:
\begin{itemize}
    \item \textbf{Routine Eye Examinations}: Comprehensive eye exams are recommended every two years for adults under 65, and annually for those over 65, individuals with a family history of glaucoma, and those diagnosed with diabetes or systemic hypertension. Early detection of asymptomatic diseases is the key to preventing permanent vision loss.
    \item \textbf{Strict Metabolic Control}: Patients with diabetes mellitus must maintain tight control of blood glucose levels (targeting an HbA1c under 7.0\% or as clinically advised), blood pressure, and cholesterol levels to minimize the risk of developing or worsening diabetic retinopathy.
    \item \textbf{Ultraviolet (UV) Protection}: Wearing sunglasses that provide 99\% to 100\% protection against UVA and UVB rays, along with a wide-brimmed hat, shields the eyes from radiation that accelerates cataractogenesis and macular degeneration.
    \item \textbf{Smoking Cessation}: Cigarette smoking is a potent, modifiable risk factor that doubles the risk of developing age-related macular degeneration and increases the likelihood of cataract formation. Cessation significantly reduces these risks over time.
    \item \textbf{Nutritional and Dietary Measures}: Consuming a balanced diet rich in leafy green vegetables (containing lutein and zeaxanthin), cold-water fish (rich in omega-3 fatty acids), and antioxidant vitamins supports macular health. For patients with intermediate AMD, taking specific high-dose vitamin formulations (based on the Age-Related Eye Disease Study 2, or AREDS2, trial) can slow progression to the advanced wet form.
    \item \textbf{Ocular Hygiene and Injury Prevention}: Practicing proper hygiene when handling contact lenses (including avoiding tap water and replacing cases regularly) prevents devastating corneal infections (microbial keratitis). Wearing certified protective eyewear (safety glasses, goggles, or face shields) during home DIY projects, gardening, sports, or industrial work prevents physical and chemical trauma.
    \item \textbf{Adoption of the 20-20-20 Rule}: To reduce digital eye strain (computer vision syndrome), individuals working on screens should look at an object at least 20 feet (6 metres) away for at least 20 seconds, every 20 minutes.
\end{itemize}

\section*{Sources and Further Reading}
\begin{itemize}
    \item World Health Organization (WHO): Blindness and Vision Impairment Fact Sheet.\\
    URL: \texttt{https://www.who.int/news-room/fact-sheets/detail/blindness-and-visual-impairment}
    \item International Agency for the Prevention of Blindness (IAPB): Vision Atlas and Global Resources.\\
    URL: \texttt{https://www.iapb.org}
    \item National Eye Institute (NEI): Eye Health Information and Clinical Research Updates.\\
    URL: \texttt{https://www.nei.nih.gov}
    \item American Academy of Ophthalmology (AAO): EyeSmart Patient Education Resources.\\
    URL: \texttt{https://www.aao.org/eye-health}
    \item Royal Australian and New Zealand College of Ophthalmologists (RANZCO): Patient Guides and Clinical Guidelines.\\
    URL: \texttt{https://ranzco.edu}
    \item Vision Australia: Assistive Technology, Low Vision Services, and Rehabilitation Support.\\
    URL: \texttt{https://www.visionaustralia.org}
\end{itemize}